%

%
\section{Empirical Evaluation}
\label{sec:experiments}
\vspace*{-3pt}

In this section, we evaluate \fsvi on high-dimensional classification tasks that were out of reach for function-space variational inference methods proposed in prior works and compare \fsvi to several well-established and state-of-the-art Bayesian deep learning and deterministic uncertainty quantification methods.
We show that \fsvi (sometimes \emph{significantly}) outperforms existing Bayesian and non-Bayesian methods in terms of their in-distribution uncertainty calibration and out-of-distribution predictive uncertainty estimation.
For a details on models, training and validation procedures, and datasets used, see~\Cref{appsec:model_details}.
For a comparison to~\citet{sun2019fbnn} on small-scale regression tasks, see~\Cref{appsec:uci}.
%



%

%
%
%
%
%
%





%



%
%
%

%
%
%


%
%
%
%
%
%
%
%



%


\vspace{-3pt}
\subsection{Predictive Performance, Uncertainty Estimation, and Distribution Shift Detection}
\label{sec:exps_setup}
\vspace*{-3pt}

In this set of experiments, we assess the reliability of the uncertainty estimates generated by \fsvi.
If a \bnn trained via \fsvi is able to perform reliable uncertainty estimation, its predictive uncertainty will be significantly higher on input points that were generated according to a different data-generating distribution than the training data.
For models trained on the FashionMNIST dataset, we use the MNIST and NotMNIST datasets as out-of-distribution evaluation points, while for models trained on the CIFAR-10 dataset, we use the SVHN dataset as out-of-distribution evaluation points.

For models trained on either FashionMNIST or CIFAR-10, we evaluate their in-distribution performance in terms of test accuracy, test log-likelihood, and test calibration.
To evaluate the quality of different models' uncertainty estimates, we compute uncertainty estimates for the pairs FashionMNIST/MNIST, FashionMNIST/NotMNIST, and CIFAR-10/SVHN to and measure for a range of thresholds how well the datasets in each pair can be separated solely based on the uncertainty estimates.
This experiment setup follows prior work by~\citet{van2020uncertainty} and~\citet{immer2020improving}.
We report the area under the receiver operating characteristic (ROC) curve in Tables~\ref{tab:results_fmnist}~and~\ref{tab:results_cifar}.
%




%

%
%
%
%
%
%
%
%
%
%
%
%
%
%
%
%
%
%
%
%
%
%
%
%
%
%
%
%
%
%
%
%



%

\setlength{\tabcolsep}{3pt}
\begin{table*}[t!]
    \caption{
        Comparison of in- and out-of-distribution performance metrics on FashionMNIST (mean $\pm$ standard error over ten random seeds).
        The last two columns show the AUROC for binary in- vs. out-of-distribution detection on MNIST (M) and NotMNIST (NM).
        MNIST and NotMNIST are used as out-of-distribution datasets.
        Best overall results for single and ensemble models are printed in boldface with gray shading.
        Results within a $95$\% confidence interval of the best overall result are printed in boldface only.
        All methods use the same four-layer CNN architecture.
        For further details about model architectures and training and evaluation protocols, see~\Cref{appsec:model_details}.
    }
    \vspace{-5pt}
    \centering
    \small
    \begin{tabular}{l c c c c}
    \toprule
    \textbf{Method}
    &  \textbf{Accuracy} $\uparrow$
    %
    &  \textbf{ECE} $\downarrow$  
    &  \textbf{AUROC} M $\uparrow$ & \textbf{AUROC} NM $\uparrow$
    %
    %
    \\
    \midrule
    \map
     & 91.73\pms{0.08}  & 0.037\pms{0.001} & 87.00\pms{0.30} & 74.85\pms{1.31}
    \\
    \mfvi~\citep{blundell2015mfvi}
     & 91.03\pms{0.04} & 0.038\pms{0.001} & 93.10\pms{0.34} & 88.88\pms{0.74}
    \\
    \mfvi (tempered)
     & 91.38\pms{0.05} & 0.058\pms{0.001} & 86.30\pms{0.29} & 80.78\pms{0.68}
    \\
    \mfvi (radial)~\citep{farquhar2020radial}
     & 90.31\pms{0.11} & 0.035\pms{0.001} & 84.40\pms{0.68} & 82.11\pms{1.15}
    \\
    \mcd~\citep{gal2016dropout}
     & 90.55\pms{0.04} & $0.012$\pms{0.001} & 88.46\pms{0.57} & 80.02\pms{1.04}
    \\
    \swag~\citep{maddox2019swag}
     & 92.56\pms{0.05} & 0.043\pms{0.001} & 85.18\pms{0.35} & 80.31\pms{0.30}
    \\
    \textsc{duq}~\citep{van2020uncertainty}
    & $92.40$\pms{0.20} & $-$ & 95.50\pms{0.70} & 94.60\pms{1.80} 
    \\
    \bnn-\textsc{laplace}~\citep{immer2020improving}
    & 92.25\pms{0.10} & $0.012\pms{0.003}$ & 95.55\pms{0.60} & ~~~~~$-$~~~~~
    \\
    \textsc{spg}~\citep{ma2021functional}
     & 91.60\pms{0.14} & ~~~~~$-$~~~~~ & 95.60\pms{6.00} & ~~~~~$-$~~~~~
    \\
    \fsvi ($p_{\bX_{\calC}}$ = random monochrome)
     & $93.13\pms{0.13}$\hspace{3pt} & $0.012\pms{0.002}$ & ${96.23}\pms{0.46}$ & ${95.02}\pms{0.69}$
    \\
    \fsvi ($p_{\bX_{\calC}}$ = KMNIST)
     & \cellcolor[gray]{0.9}$\mathbf{93.48}\pms{0.12}$\hspace{3pt} & \cellcolor[gray]{0.9}$\mathbf{0.010}$\pms{0.001} & \cellcolor[gray]{0.9}$\mathbf{99.80}\pms{0.20}$ & \cellcolor[gray]{0.9}$\mathbf{97.26}\pms{0.23}$
    %
    %
    %
    \\[1pt]
    \cline{1-5}
    \\[-7pt]
    Deep Ensemble
    %
     & 92.49\pms{0.01} & \cellcolor[gray]{0.9}$\mathbf{0.019}$\pms{0.000}\hspace{3pt} & 89.22\pms{0.09} & 83.17\pms{0.91}
    %
    %
    %
    \\
    \fsvi Ensemble ($p_{\bX_{\calC}}$ = random monochrome)
     & \cellcolor[gray]{0.9}$\mathbf{94.44}$\pms{0.07}\hspace{3pt} & 0.020\pms{0.001} & \cellcolor[gray]{0.9}$\mathbf{97.85}$\pms{0.15} & \cellcolor[gray]{0.9}$\mathbf{96.95}\pms{0.20}$
    \\
    \bottomrule
    \end{tabular}
    \label{tab:results_fmnist}
    %
\end{table*}

%
%
%
%
%
%
%
%
%
%
%
%
%
%
%
%
%
%
%
%
%
%
%
%
%
%
%
%
%
%
%
%
%
%
%
%
%
%
%
%
%
%
%
%
%
%
%
%
%
%
%
%
%
%
%
%
%
%
%
%
%
%
%
%
%
%
%
%
%
%
%
%
%
%
%
%
%
%

%


\textbf{Predictive Performance and Calibration.}$~$
%
%
%
%
%
To assess in-distribution predictive performance and calibration, we report the test accuracy, negative log-likelihood (NLL), and expected calibration error (ECE) for models trained on FashionMNIST and CIFAR-10 in Tables~\ref{tab:results_fmnist}~and~\ref{tab:results_cifar}.
On both \mbox{FashionMNIST} and CIFAR-10, \fsvi achieves the lowest NLL and either the best or second-best predictive accuracy and ECE, respectively, across all methods.
Notably, \fsvi significantly outperforms \textsc{spg}~\citep{ma2021functional}, an alternative function-space variational inference method.
%
%
%
%


%


\textbf{Predictive Uncertainty under Distribution Shift.}$~$
In Tables~\ref{tab:results_fmnist}~and~\ref{tab:results_cifar}, we report evaluation metrics that elucidate the reliability of different methods' predictive uncertainty under distribution shift.
\fsvi exhibits reliable predictive uncertainty estimates that allow distinguishing between in- and out-of-distribution inputs with high accuracy.
%
As would be expected, we observe that using context distributions that reflect our knowledge about the data-generating process can significantly improve uncertainty quantification under \fsvi.
For the FashionMNIST experiment, we used the KMNIST dataset, which contains grayscale images of Kuzushiji letters, and for the CIFAR-10 experiment, we used the CIFAR-100 dataset, which contains RGB images of 100 classes.
Both KMNIST and CIFAR-100 differ from the OOD datasets (MNIST and NotMNIST and SVHN, respectively) used to compute OOD-AUROC metrics in Tables~\ref{tab:results_fmnist} and~\ref{tab:results_cifar}, but using them as context distributions significantly increased the ability of \bnns trained via \fsvi to identify distributionally shifted samples.
Since the variational objective encourages matching the prior (which we chose to have high variance) on samples from the context distribution can improve uncertainty estimation in regions of the input space far from the training data.




%

\setlength{\tabcolsep}{2pt}
\begin{table*}[t!]
    \caption{
        Comparison of in- and out-of-distribution performance metrics on CIFAR-10 (mean $\pm$ standard error over ten random seeds).
        SVHN and corrupted CIFAR-10 (C-CIFAR) are used as an out-of-distribution datasets.
        The penultimate column shows the AUROC for binary in- vs. out-of-distribution detection on SVHN.
        Best overall results for single and ensemble models are printed in boldface with gray shading.
        Results within a $95$\% confidence interval of the best overall result are printed in boldface only.
        All methods use a ResNet-18 architecture.
        For further details about model architectures and training and evaluation protocols, see~\Cref{appsec:model_details}.
    }
    \vspace{-5pt}
    \centering
    \small
    \begin{tabular}{l c c c c}
    \toprule
    \textbf{Method}
    &  \textbf{Accuracy}$\uparrow$
    %
    &  \textbf{ECE}$\downarrow$  
    &  \textbf{OOD-AUROC}$\uparrow$ & \textbf{C-CIFAR Acc}$\uparrow$
    %
    %
    \\
    \midrule
    \\[-9pt]
    \map
     & 93.19\pms{0.11} & 0.043\pms{0.001} & 94.65\pms{0.27} & 78.87\pms{1.39}
    \\
    \mfvi~\citep{blundell2015mfvi}
     & 89.98\pms{0.09} & 0.040\pms{0.001} & 92.14\pms{0.34} & 79.36\pms{1.35}
    \\
    \mfvi (tempered)
     & 90.87\pms{0.11} & 0.048\pms{0.001} & 91.82\pms{0.90} & 79.86\pms{1.32}
    \\
    \mcd~\citep{gal2016dropout}
     & 93.55\pms{0.07} & 0.040\pms{0.001} & 92.44\pms{0.57} & 80.13\pms{1.37}
    \\
    \swag~\citep{maddox2019swag}
     & 93.13\pms{0.14} & 0.067\pms{0.002} & 89.79\pms{0.50} & 76.12\pms{0.51}
    \\
    \textsc{vogn}~\citep{osawa2019practical}
    & 84.27\pms{0.20} & 0.040\pms{0.002} & 87.60\pms{0.20} & ~~~~~$-$~~~~~
    \\
    \textsc{duq}~\citep{van2020uncertainty}
    & \cellcolor[gray]{0.9}$\mathbf{94.10}$\pms{0.20}\hspace{3pt} & $-$ & 92.70\pms{1.30} & ~~~~~$-$~~~~~
    \\
    \textsc{spg}~\citep{ma2021functional}
     & 77.69\pms{0.64} & ~~~~~$-$~~~~~ & 88.30\pms{4.00} & ~~~~~$-$~~~~~
    \\
    \fsvi ($p_{\bX_{\calC}}$ = random monochrome)
     & $93.35$\pms{0.04} & 0.034\pms{0.001} & $94.76$\pms{0.24} & $\mathbf{80.81}$\pms{0.43}
    \\
    \fsvi ($p_{\bX_{\calC}}$ = CIFAR-100)
     & $93.57$\pms{0.04} & \cellcolor[gray]{0.9}$\mathbf{0.026}$\pms{0.001} & \cellcolor[gray]{0.9}$\mathbf{98.07}$\pms{0.10} & \cellcolor[gray]{0.9}$\mathbf{81.20}$\pms{0.42}
    %
    %
    %
    %
    \\[1pt]
    \cline{1-5}
    \\[-7pt]
    Deep Ensemble
     & $95.13$\pms{0.06} & 0.019\pms{0.001} & $\mathbf{98.04}$\pms{0.07} & $\mathbf{81.22}$\pms{0.37}
    \\
    \fsvi Ensemble ($p_{\bX_{\calC}}$ = random monochrome)
     & \cellcolor[gray]{0.9}$\mathbf{95.19}$\pms{0.03}\hspace{3pt} & \cellcolor[gray]{0.9}\textbf{0.013}\pms{0.001} & \cellcolor[gray]{0.9}$\mathbf{99.19}$\pms{0.41} & \cellcolor[gray]{0.9}$\mathbf{81.35}$\pms{0.48}
    \\
    \bottomrule
    \end{tabular}
    \label{tab:results_cifar}
    %
\end{table*}

%
%
%
%
%
%
%
%
%
%
%
%
%
%
%
%
%
%
%
%
%
%
%
%
%
%
%
%
%
%
%
%
%
%
%
%
%
%
%
%
%
%
%
%
%
%
%
%
%
%
%
%
%
%
%
%
%
%
%
%
%
%
%
%
%
%
%
%
%
%
%
%
%


\vspace*{-3pt}
\subsection{Generalization and Reliability of Predictive Uncertainty under Distribution Shift}
\vspace*{-3pt}

To assess the reliability of predictive models in deep learning,~\citet{ovadia2019uncertainty} propose the following desiderata:
In order for a model to be considered reliable, it ought to (i) exhibit low predictive uncertainty on training data and high predictive uncertainty on out-of-distribution inputs, (ii) generate predictive uncertainty estimates that allow distinguishing in- from out-of-distribution inputs, and (iii) if possible, maintain high predictive accuracy even under distribution shift.
Models that satisfy these desiderata are less likely to make poor, high-confidence predictions and more amenable for use in safety-critical downstream tasks.

%
\begin{wrapfigure}{r}{0.45\textwidth}
\centering
\vspace*{-26pt}
    \includegraphics[width=0.85\linewidth, trim={5pt 0 5pt 0},clip]{figures/neurips2021/mnist/rotated_mnist}
    \vspace{-5pt}
    \caption{
        Predictive uncertainty and accuracy on rotated MNIST.
        Models with reliable uncertainty estimates would exhibit higher predictive uncertainty the more the digits are rotated.
        Ideally, such models would maintain high predictive accuracy (low Brier score).
    }
    \label{fig:rotated_mnist}
\vspace*{-19pt}
\end{wrapfigure}
%

To illustrate these desiderata, we follow~\citet{ovadia2019uncertainty} and consider the rotated MNIST task, where a model is trained on MNIST and evaluated on rotated MNIST digits.
The goal is to maintain a high level of predictive accuracy (measured in terms of Brier scores) while exhibiting an increasing level of predictive uncertainty on distribution shifts of increasing magnitude.
\Cref{fig:rotated_mnist} shows Brier scores (lower is better) and predictive entropy estimates (higher means more uncertain) of four different models.
As rotating the MNIST digits gradually shifts the data distributions, we would expect Brier scores to increase (corresponding to worse predictive accuracy) as the rotation angle increases.
A model with reliable predictive entropy estimates would only experience a small decrease under distribution shift while exhibiting a large increase in predictive uncertainty.
As can be seen in the plot, the Brier scores of \fsvi decreases the least, while \fsvi's uncertainty is significantly higher than other models'.
To assess the reliability of different uncertainty quantification methods on a more challenging distribution-shift task, we consider corrupted CIFAR-10 inputs under the second-mildest corruption level used in~\citep{ovadia2019uncertainty} and report our results in~\Cref{tab:results_cifar}.
Consistent with the rotated MNIST results, \fsvi achieves the highest accuracy on the corrupted data.


%

\vspace*{-3pt}
\subsection{Safety-Critical Uncertainty-Aware Selective Prediction: Diabetic Retinopathy Diagnosis}
\label{sec:diabetic}
\vspace*{-3pt}

%
\begin{wrapfigure}{r}{0.25\textwidth}
\centering
\vspace*{-15pt}
    \hspace*{-15pt}\includegraphics[width=\linewidth]{figures/uai2022/retina/healthy_3.pdf}
    \\
    \hspace*{-15pt}\includegraphics[width=\linewidth]{figures/uai2022/retina/unhealthy_3.pdf}
    \vspace*{-5pt}
    \caption{
        Retina scan examples. \textbf{Top:} healthy. \textbf{Bottom:} unhealthy.
    }
    \label{fig:retina}
\vspace*{-10pt}
\end{wrapfigure}
%

To evaluate the reliability of the predictive uncertainty of \fsvi in a real-world safety-critical setting, we consider the task of diagnosing diabetic retinopathy (DR), a medical condition that can lead to impaired vision, from retina scans~\citep{leibig2017leveraging,filos2019systematic,Band2021benchmarking}.
We use two publicly available datasets,~\citet{kaggle_2015} and~\citet{APTOS_2019}, each containing RGB images of a human retina graded by a medical expert on the following scale: 0 (no DR), 1 (mild DR), 2 (moderate DR), 3 (severe DR), and 4 (proliferative DR).
The Kaggle dataset was collected from patients in the United States, while the APTOS dataset was collected from patients in India using cheaper but more modern scanning devices.
We follow~\citet{leibig2017leveraging},~\citet{filos2019systematic}, and~\citet{Band2021benchmarking} and binarize all examples from both the EyePACS and APTOS datasets by dividing the classes up into sight-threatening diabetic retinopathy---defined as moderate diabetic retinopathy or worse (classes $\{2,3,4\}$)---and non-sight-threatening diabetic retinopathy---defined as no or mild diabetic retinopathy (classes $\{0,1\}$).
This results in a binary prediction task.

To assess the reliability of predictive models when medical training and test data are obtained from different patient populations or collected with the same medical equipment, we follow~\citet{Band2021benchmarking} and use the Kaggle dataset for training and the distributionally shifted APTOS dataset for evaluation.
The results are shown in~\Cref{fig:country_shift}, which plot the ROC curves for the binary prediction problems as well as the area under the ROC curve for an uncertainty aware selective prediction task.
For further details about the uncertainty-aware selective prediction evaluation protocol, see~\Cref{appsec:retina}.
\Cref{fig:country_shift} shows that \fsvi performs well on all four tasks and is only outperformed by \mcd.
For full tabular results, see~\Cref{appsec:retina_tabular}.








\begin{figure*}[t!]
\centering
\captionsetup[subfigure]{justification=centering}
\includegraphics[width=\linewidth]{figures/uai2022/retina/legend.pdf}\\[-10pt]
\hspace*{-10pt}\subfloat[
    ROC: In-Domain
]{\includegraphics[width=0.24\linewidth]{figures/uai2022/retina/roc-aptos-in_domain_test-indomain-drd.pdf}}%
\subfloat[
    ROC: Country Shift
]{\includegraphics[width=0.24\linewidth]{figures/uai2022/retina/roc-aptos-ood_test-indomain-drd.pdf}}%
\subfloat[
    Selective Prediction AUROC: In-Domain
]{\includegraphics[width=0.24\linewidth]{figures/uai2022/retina/retention-aptos-in_domain_test-indomain-AUC-no_oracle.pdf}}%
\subfloat[
    Selective Prediction AUROC: Country Shift
]{\includegraphics[width=0.24\linewidth]{figures/uai2022/retina/retention-aptos-ood_test-indomain-AUC-no_oracle.pdf}}%
\vspace*{-3pt}
\caption{
%
    We jointly assess model predictive performance and uncertainty quantification on both in-domain and distributionally shifted data.
    \textbf{Left:} The \textit{receiver operating characteristic curve} (ROC) for in-population diagnosis on the (\textbf{a})~\citet{kaggle_2015} test set and for (\textbf{b}) changing medical equipment and patient populations on the~\citet{APTOS_2019} test set. 
    The dot in 
    \textbf{black}
    %
    denotes the NHS-recommended 85\% sensitivity and 80\% specificity ratios~\citep{widdowson2016management}.
    \textbf{Right:} \emph{Selective prediction} on AUROC in (\textbf{c})~\citet{kaggle_2015} and (\textbf{d})~\citet{APTOS_2019} settings.
    Shading denotes standard error over six random seeds.
    See~\Cref{appsec:retina_tabular} for tabular results.
}
\label{fig:country_shift}
\vspace*{-10pt}
\end{figure*}









%
%

%
%
%
%
%
%
%
%
%
%
%
%
%
%
%
%
%
%
%
%
%
%

%

%
%
%
%
%
%
%
%
%
%
%
%
%
%
%

%

%
%
%
%
%







%



%
%
%

%
%
%
%




%

%
%
%
%
%
%
%
%
%
%
%
%
%
%
%
%
%
%
%
%
%
%
%
%
%
%

%
